% Study 3, the registered design as the dated addition of 2026-10-16, item 7, fixes it. Input from
% paper.tex; every number traces to sources.md and is asserted by check_numbers.py. The remaining detail
% (the dynamic law, the fallback, completeness and the source-gap rule, the fallback domain, the gate, the
% seeds, the contamination ledger) is in study3-detail.tex (Appendix I). Macros are provided in case the
% file is compiled before paper.tex defines them.
\providecommand{\qs}{QScore9}
\providecommand{\hc}{\texttt{hurdle\_count\_28}}

\section{Study 3: the registered design}
\label{sec:s3design}

\paragraph{Domain, frame and strata.} The series key is the problem type by community board, enumerated
from the frozen 2024 inventory of 9,496 cells captured on 2026-09-14, of which 8,469 carry a recognized
board and problem and form the whole eligible frame, with no candidate pool. Counts are requests by
creation date in New York calendar days with daylight saving time, deduplicated by request identifier with
the latest captured revision kept, counting requests and not closures. The strata are Study 2's, computed
on the context zero fraction of 2026-03-05 through 2026-06-02: $[0, 0.5)$ dense, $[0.5, 0.9)$ sparse,
$[0.9, 1]$ very sparse; with selection seed 20261001 the first 100 complete series of each stratum in
selection-hash order are taken, with no transfer of slots. The context from 2025-09-01 through 2026-06-02
is inspected for completeness and the strata only. For this event-enumerated source a day is complete when
its capture is exhausted, and a calendar day with no request in any frame cell is a source gap that stops
the domain before any scoring, a rule added because npm's served zero days were found only after scoring.
PyPI through ClickPy was the sole registered fallback, for an operational failure of NYC 311 by 2026-10-27
alone, and was not invoked.

\paragraph{Calendar, models and variants.} The calendar is Study 2's (Section~\ref{sec:s2design}): context
from 2025-09-01, cutoffs 2026-06-03, 2026-07-03 and 2026-08-02 with 275, 305 and 335 context days,
horizons 7 and 30, endpoint 2026-08-31, in New York calendar days. The three foundation models are the
Study 1 pins, with Study 2's completion, scores and baselines unchanged, and all 22 registered variants are
forecast at every origin. For each model four variants are scored from the same saved base forecast with
no refitting: \emph{completed}, the control; \emph{hurdle\_fixed}, Study 2's hurdle with the 28-day
Jeffreys occurrence; \emph{hurdle\_dynamic}, new, which takes at each lead the occurrence probability of the
registered \texttt{dynamic\_count} law, a discounted Beta-Bernoulli with discount 0.98 started at $a_B =
b_B = 0.5$, with the model's completed positive law; and \emph{occurrence\_model\_size\_count}, new, the
fourth factorial cell, the model's own completed occurrence $1 - F(0)$ with \hc{}'s positive law. Where
\texttt{dynamic\_count} falls back to \hc{}'s law, with fewer than two request days in the last 28 of a
context, \emph{hurdle\_dynamic} equals \emph{hurdle\_fixed} for that window. No variant is fitted or
selected on targets.

\paragraph{The confirmatory family.} For each stratum, model and horizon, T1 compares the model's
\emph{hurdle\_dynamic} with its completed control and T2 compares it with \texttt{hurdle\_empirical\_28}:
$3 \times 3 \times 2 \times 2 = 36$ \qs{} tests, R01 to R36, listed in Table~\ref{tab:s3structure} and
committed before the gate. Differences are variant minus reference, so negative favors the variant. H3
predicts a positive T1 median in the dense stratum and a negative one in the sparse and very sparse strata,
and nothing for T2; the tests are two-sided as registered, with the direction read from the median and
reported beside the mean. Within each stratum origins are averaged within series and horizon, the 100 series
weigh equally and are the pairing and bootstrap units (20,000 replicates, pointwise 95 percent percentile
intervals, bootstrap seed 20261004), the two-sided approximate Wilcoxon signed-rank test omits zero
differences and reports $p = 1$ when every difference is zero, and Benjamini-Hochberg is applied once at
$q = 0.05$ over the 36. A stratum with fewer than ten complete series would have made its twelve tests
unavailable, with the family not shrunk by any outcome; none was. Everything else, the fixed hurdle and the
fourth cell against the control and the count hurdles per stratum, the pooled 300-series tables in Study
2's form, the scaled score, the Brier score, the CRPS, the PIT, the accounting and the concentration
descriptors, is declared secondary and adds no test. The gate, the seeds, the smoke and the contamination
ledger are in Appendix~\ref{app:s3detail}.

\begin{table}[t]
\centering\small
\caption{The confirmatory family of Study 3, R01 to R36, in its registered order: for each stratum of 100
series, each model and each horizon, T1 is the model's \emph{hurdle\_dynamic} against its own completed
control and T2 is \emph{hurdle\_dynamic} against \texttt{hurdle\_empirical\_28}. One Benjamini-Hochberg
adjustment at $q = 0.05$ covers the 36.}
\label{tab:s3structure}
\begin{tabular}{llcccc}
\toprule
 & & \multicolumn{2}{c}{$h = 7$} & \multicolumn{2}{c}{$h = 30$} \\
\cmidrule(lr){3-4} \cmidrule(lr){5-6}
Stratum & Model & T1 & T2 & T1 & T2 \\
\midrule
dense & Chronos-2 & R01 & R02 & R03 & R04 \\
dense & TimesFM 3 & R05 & R06 & R07 & R08 \\
dense & TimesFM 2.5 & R09 & R10 & R11 & R12 \\
sparse & Chronos-2 & R13 & R14 & R15 & R16 \\
sparse & TimesFM 3 & R17 & R18 & R19 & R20 \\
sparse & TimesFM 2.5 & R21 & R22 & R23 & R24 \\
very sparse & Chronos-2 & R25 & R26 & R27 & R28 \\
very sparse & TimesFM 3 & R29 & R30 & R31 & R32 \\
very sparse & TimesFM 2.5 & R33 & R34 & R35 & R36 \\
\bottomrule
\end{tabular}
\end{table}
